The empirical transition operator defines a stable finite-state baseline, but the two-step operator differs from the Markov prediction in structured ways. The residual \(\Delta\) is not uniformly diffuse: heatmaps show localized positive and negative transition blocks, and the singular spectrum concentrates energy in the first few modes. In the computed range, approximately three to five modes capture most of the residual energy. Low-rank corrections reduce two-step prediction error, indicating that the leading residual modes encode reproducible transition information rather than only noise.

Control comparisons strengthen this interpretation. Iid and residue-balanced shuffles substantially reduce residual strength, singular values, and mutual-information excess. Markov synthetic sequences preserve first-order statistics but do not reproduce the full two-step residual structure. Block and window-preserving resampling retain more of the structure, suggesting that the signal is related to local sequential constraints rather than marginal residue counts alone.
